\section{第三人}
\begin{enumerate}
  \item 有独三。\zd{一般是主张权利的}。在记之前，想想老师给你讲的原理。
        \begin{enumerate}
          \item 跟原被告基于一个\zd{独立的法律关系}。
          \item 程序。
                \begin{enumerate}
                  \item \zd{有独三之诉}。自己起诉，法院不能追加。
                  \item 法庭辩论终结前。
                  \item 二审只准调解不准判决，调解不成\zd{撤销原判，发回重审}（另行起诉会导致\zd{判决冲突}）。
                \end{enumerate}
        \end{enumerate}

  \item 无独三。没有独立的请求权，但与\zd{案件处理结果}有法律上的\zd{利害关系}。
        \begin{enumerate}
          \item 自己申请 \hz 法院追加。
          \item 诉讼权利。
                \begin{enumerate}
                  \item 绝对不能享有的权利。不能提\zd{管辖权异议}；不能放弃、变更诉讼请求；不能撤诉。
                  \item 附条件（\zd{是否承担义务、责任}）下的权利。上诉、签收调解书。
                \end{enumerate}
        \end{enumerate}

        \begin{kk}{有独三跟共同原告、无独三跟共同被告区别}
          先找原被告 \yjt \zd{原被告间的法律关系（诉讼标的）} \yjt 把第三人拿进去套 \yjt 套得进告知追加共同原告/共同被告 \yjt 套不进 \yjt 有独三还是无独三。
          \zd{理论基础。共同原告、共同被告前提他得是适格原告、被告}。
          \tcblower
          \al{李独资设立甲公司，李将全部资产投入其妻项目中，乙对甲公司有 200W 贷款到期，乙起诉李。\\
            A. 变更甲为被告 ×。 B. 甲列为第三人 ×。 C. 追加甲为共同被告 ×。 D. 告知追加甲为共同被告。}
        \end{kk}
\end{enumerate}
